% paper/sections/availability.tex — Availability and Installation section
% Filled in Phase 88 Plan 05 (MANU-07).
\section{Availability and Installation}
\label{sec:availability}

This paper describes \texttt{fdars}~\FdarsVersion{}.  Version~\FdarsVersion{},
described here, is available from the Python Package Index (PyPI) and can be
installed with:
\begin{verbatim}
pip install fdars==0.13.0
\end{verbatim}
Prebuilt binary wheels are provided for Linux (x86\_64 and aarch64), macOS
(x86\_64 and aarch64), and Windows (x86\_64).
The package targets the CPython stable ABI (\texttt{abi3-py39}) and supports
Python~3.9 through 3.14.

An optional \texttt{plot} extra installs \texttt{matplotlib} for the built-in
visualisation utilities:
\begin{verbatim}
pip install "fdars[plot]"
\end{verbatim}
An optional \texttt{sklearn} extra installs \texttt{scikit-learn} and enables
the \texttt{fdars.sklearn} estimator layer, which exposes \nsklearnestimators{}
estimators compatible with the scikit-learn API:
\begin{verbatim}
pip install "fdars[sklearn]"
\end{verbatim}

The advisor and MCP server are installed with the \texttt{advisor} and
\texttt{mcp} extras; provider-specific extras (\texttt{openai},
\texttt{gemini}, \texttt{ollama}) select the language-model backend (the
\texttt{mcp} and \texttt{gemini} extras require Python~3.10 or later).
The paper's reproduction scripts, reference data, and agreement checks live in
the \texttt{paper/} directory of the repository, and a \texttt{CITATION.cff}
file gives the preferred citation.

The source code is released under the MIT licence and is hosted at
\url{https://github.com/sipemu/pyfda}.
Documentation, including worked examples and API reference, is available at
\url{https://sipemu.github.io/pyfda/}.
